\clearpage
\fancyhead[L]{\small ANWENDUNG DER WAHRSCHEINLICHKEITSRECHNUNG}
\pdfbookmark[0]{Nachlass. Anwendung der Wahrscheinlichkeitsrechnung auf die Bestimmung der Bilanz für Witwenkassen}{witwen}
\SourceLeaf{124}
\begin{center}
{\Large NACHLASS.\par}\bigskip
{\large [ANWENDUNG DER WAHRSCHEINLICHKEITSRECHNUNG\par}
\medskip{\Large AUF DIE BESTIMMUNG DER BILANZ\par}\medskip
{\large FÜR WITWENKASSEN ]\par}
\medskip\rule{.15\linewidth}{.3pt}\par\bigskip
{\large [I.]\par}
\pdfbookmark[1]{I. Allgemeine Uebersicht der Methode}{witwen-I}
\textit{[Allgemeine Uebersicht der Methode.]}\par
\textit{[Auszug aus einem Votum bei der schriftlichen Abstimmung im Universitäts-Senat.]}\par
\end{center}
Das vorstehende [hier eingeklammerte] Votum des Herrn Prof. D.: [Das Königl. Univ. Curatorium scheint zu befürchten, dass bei der grossen Anzahl der jetzt vorhandenen Witwen die Kasse über lang oder kurz nicht im Stande sein werde, die jetzt auf 250 Thl. angewachsenen Pensionen zu bestreiten. Es verlangt daher einen Bericht darüber, ob gegründete Ursache zu einer solchen Besorgniss vorhanden sei, und durch welche Mittel die etwa drohende Gefahr abgewendet werden könne. .\,.\,.\,] spricht den eigentlichen \textit{Fragepunkt} so treffend aus, dass ich der ersten Hälfte dieses Votum nur wörtlich beitreten kann. Wenn in Zweifel gezogen ist, ob die Kasse im Stande sein werde, der ihr obliegenden Verpflichtung nachhaltig zu genügen, so ist dies doch wahrlich der ungeeignetste Zeitpunkt, \textit{grössere} Anforderungen an die Kasse zu stellen.

\indent Ich kann mich der öffentlichen Meinung über diese Anstalt noch bis 40 Jahr rückwärts erinnern. Damals schon galt sie für ein herrliches Kleinod der Universität, einzig in seiner Art, und zwar gerade wegen ihrer Eigenthümlichkeiten. \textit{Vollkommene} Freiheit, ob man beitreten wolle oder nicht, ja, mit einer vergleichungsweise geringen Aufopferung, wieder einzutreten, wenn man ausgetreten war; ein sehr geringer Beitrag. Und damals betrug die Pension nur 150 oder 160 Thl. Nicht die Grösse der Pension war das Anziehende, sondern die \textit{liberale Art, wie} dem, der Göttingischer Professor werden konnte, eine sichere Unterstützung einer nachbleibenden Witwe, mit der Aussicht, unter der weisen Verwaltung sie nach und nach noch erhöhet zu erhalten, dargeboten wurde. Wer mehr wünschte, betheiligte sich noch nebenbei in einer andern Witwenkasse. Jetzt ist nun die Pension auf 250 Thl. gestiegen, und die liberale Art ist bis heute dieselbe geblieben. Gebe Gott, dass niemals nöthig werde, an dieser Art \textit{irgend etwas} zu ändern! Zwangsprocente auf das Gehalt, um durch Drehung am Stundenzeiger das zu erhalten, was nur der allmählige Fortschritt des Minutenzeigers gewähren kann, würde nicht blos viel zu unwirksam sein um den Zweck zu erreichen, sondern den noblen Charakter der Anstalt ganz zerstören.
